\section{Problem 1}

\subsection{} %1.1

Since the solid cone $\cQ^3$ is closed, convex,
and rotational symmetric along the first axis ($t$ axis),
we only need consider a small set $S$ of candidate points
when trying to find the projection of a point $y$ into this set $\cQ^3$.
This is because the point projected to must be unique
and thus we can eliminate points that have `brother/sister' points which
shares the same distance to $y$ due to symmetry.
Specifically, we can choose $S$ be the (or a) plane that contains
both the point $y$ and the first axis.
The problem then reduces to a simple two-dimensional geometry problem
of middle school level (Figure~\ref{fig:1}).
We can directly write down the solution as specified in the question.

\begin{figure}[ht]
\centering
    \includegraphics[width=0.36\linewidth]{figure/fig1}
    \caption{\small
    The problem of projecting point $y=(t,u)$ to $\cQ^3$ is
    reduced to this two dimensional problem.}
    \label{fig:1}
\end{figure}

\subsection{} %1.2

Let point $z=\Pi_H(y)=(a,b,a)\in H$ be the projected point.
Then $z-y=(a-s, b-u_1, a-u_2)$ must parallel to the normal direction of plane $H$,
namely $n=(1,0,-1)$.
Solve this constraint,
we get $b=u_1$ and $a=(s+u_2)/2$.

\subsection{} %1.3

Let $p=(t,u_1,u_2)\in \cQ^3\cap H$ be a point in the intersection.
Since $p\in H$, we have $t=u_2$.
Since $p \in \cQ$, we have $u_1^2+u_2^2=u_1^2 + t^2\le t^2$,
and thus $u_1\equiv0$.
Together with $t\ge\|x\|_2\ge0$,
we know the intersection is contained in the ray $R:=\{t(1,0,1): t\ge0\}$.
It's easy to see that this ray also contained in the intersection.
So $\cQ^3\cap H=R$.
Then the projection of point $y$ to this ray $R$ is a simple optimization problem
\begin{equation}
    t^*=\argmin_{t\ge0}\|(s,u_1,u_2)-t(1,0,1)\|_2^2=\argmin_{t\ge0}(s-t)^2+(u_2-t)^2.
\end{equation}
The solution is $t^*=(s+u_2)_+/2$.
So comes the desired answer.

\subsection{} %1.4

Use previous results,
we have
\begin{equation}
\begin{split}
    \dist(x^k, \cQ^3)
        &= \bigg\|(k,1,k)-\frac{k+\sqrt{1+k^2}}{2\sqrt{1+k^2}}(\sqrt{1+k^2},1,k)\bigg\|_2 \\
        &= \bigg\|\bigg( \frac{\sqrt{1+k^2}-k}{2},\frac{k}{2\sqrt{1+k^2}}-\frac{1}{2},\frac{k}{2\sqrt{1+k^2}}\big(k-\sqrt{1+k^2}\big) \bigg)\bigg\|_2 \\
        &\rightarrow 0\quad(k\rightarrow\infty),
\end{split}
\end{equation}
\begin{equation}
    \dist(x^k, H) = 0,
\end{equation}
and
\begin{equation}
    \dist(x^k, \cQ^3\cap H) = 1.
\end{equation}
It is evident that the claim in the question holds.
